\documentclass{article}%
\usepackage{amsmath}
\usepackage{amsfonts}
\usepackage{amssymb}
\usepackage{graphicx}%
\begin{document}
\section{\bigskip Correspondance between the DKE\ code and the Delphine
absorption code}

\subsection{Derivation of the quasilinear diffusion coefficient}

We want to show how the equations of the absorption code can be derived from
the more general, non-circular, bounce-averaged equations solved by DKE. The
quasilinear operator in DKE is given by [1]-()
\begin{equation}
\overline{D}_{b,n}^{\text{RF}}(p,\xi)=\overline{D}_{b,n,0}^{\text{RF}%
}\left\vert \Theta_{\mathbf{k}}^{b,\left(  n\right)  }\right\vert ^{2}%
\frac{\gamma p_{Te}}{p\left\vert \xi\right\vert }\delta\left(  N_{b\Vert
}-N_{\Vert\text{res}}\right)  2\pi\delta\left(  \theta-\theta_{b}\right)
\label{difcoeffnorm2}%
\end{equation}
with
\begin{align}
\overline{D}_{b,n,0}^{\text{RF}}  &  =\frac{1}{\nu_{e}p_{Te}^{2}}\frac
{1}{v_{Te}}\frac{\pi}{\omega_{b}}\frac{e^{2}}{\epsilon_{0}}\frac{1}{4\pi
^{2}rR}\left\vert \widehat{\psi}\cdot\widehat{r}\right\vert \frac
{f_{\text{inc},b}^{l+1/2}}{\left\vert \mathbf{\Phi}_{b}\right\vert
}P_{b,\text{inc}}\nonumber\\
&  =\frac{1}{rR}\frac{1}{m_{e}\ln\Lambda}\frac{1}{\omega_{b}\omega_{pe}^{2}%
}\left\vert \widehat{\psi}\cdot\widehat{r}\right\vert \frac{f_{\text{inc}%
,b}^{l+1/2}}{\left\vert \mathbf{\Phi}_{b}\right\vert }P_{b,\text{inc}}
\label{drf0}%
\end{align}
where $\omega_{pe}$ is the electron plasma frequency
\begin{equation}
\omega_{pe}=\sqrt{\frac{e^{2}n_{e}}{\epsilon_{0}m_{e}}}%
\end{equation}


We also recall the expressions (\ref{thetakn}), (\ref{nparres}) and
(\ref{phib2}-\ref{phibp2})%
\begin{equation}
\Theta_{\mathbf{k}}^{b,\left(  n\right)  }=\frac{1}{\sqrt{2}}e_{b0,+}%
e^{-i\alpha_{b}}J_{n-1}\left(  \frac{k_{b\bot}v_{\bot}}{\Omega}\right)
+\frac{1}{\sqrt{2}}e_{b0,-}e^{+i\alpha_{b}}J_{n+1}\left(  \frac{k_{b\bot
}v_{\bot}}{\Omega}\right)  +\frac{p_{\Vert}}{p_{\bot}}e_{b0,\Vert}J_{n}\left(
\frac{k_{b\bot}v_{\bot}}{\Omega}\right)  \label{thetakn2}%
\end{equation}%
\begin{equation}
N_{\Vert\text{res}}=\frac{c}{v_{\Vert}}\left(  1-\frac{n\Omega}{\omega_{b}%
}\right)  =\frac{1}{\beta_{Te}}\frac{p_{Te}}{p_{\Vert}}\left(  \gamma
-\frac{n^{\prime}\omega_{ce}}{\omega_{b}}\right)  \label{nparres2}%
\end{equation}%
\begin{align}
\mathbf{\Phi}_{b}  &  =\mathbf{\Phi}_{bP}+\mathbf{\Phi}_{bT}\label{phib3}\\
\mathbf{\Phi}_{bP}  &  =\text{Re}\left[  \mathbf{N}_{b}-\left(  \mathbf{N}%
_{b}\cdot\mathbf{e}_{b}\right)  \mathbf{e}_{b}^{\ast}\right] \label{phibp3}\\
\mathbf{\Phi}_{bT}  &  =\mathbf{-}\frac{1}{2}\mathbf{e}_{b}^{\ast}\cdot
\frac{\partial\mathbb{{K}^{H}}}{\partial\mathbf{N}}\cdot\mathbf{e}_{b}
\label{phibt3}%
\end{align}


For the lower-hybrid wave, the wave-particle resonance occurs through Landau
damping, $n=0$, and the diffusion coefficient becomes%
\begin{equation}
\overline{D}_{b,0}^{\text{RF}}(p,\xi)=\overline{D}_{b,0,0}^{\text{RF}%
}\left\vert \Theta_{\mathbf{k}}^{b,\left(  0\right)  }\right\vert ^{2}%
\frac{\gamma p_{Te}}{p\left\vert \xi\right\vert }\delta\left(  N_{b\Vert
}-N_{\Vert\text{res}}\right)  2\pi\delta\left(  \theta-\theta_{b}\right)
\end{equation}
with
\begin{equation}
\Theta_{\mathbf{k}}^{b,\left(  0\right)  }=\frac{1}{\sqrt{2}}e_{b0,+}%
e^{-i\alpha_{b}}J_{-1}\left(  \frac{k_{b\bot}v_{\bot}}{\Omega}\right)
+\frac{1}{\sqrt{2}}e_{b0,-}e^{+i\alpha_{b}}J_{1}\left(  \frac{k_{b\bot}%
v_{\bot}}{\Omega}\right)  +\frac{p_{\Vert}}{p_{\bot}}e_{b0,\Vert}J_{0}\left(
\frac{k_{b\bot}v_{\bot}}{\Omega}\right)
\end{equation}
and%
\begin{equation}
N_{\Vert\text{res}}=\frac{\gamma}{\beta_{Te}}\frac{p_{Te}}{p_{\Vert}}%
\end{equation}


Using the Bessel function relation () we find that%
\begin{equation}
\Theta_{\mathbf{k}}^{b,\left(  0\right)  }=\frac{1}{\sqrt{2}}\left[
e_{b0,-}e^{+i\alpha_{b}}-e_{b0,+}e^{-i\alpha_{b}}\right]  J_{1}\left(
\frac{k_{b\bot}v_{\bot}}{\Omega}\right)  +\frac{p_{\Vert}}{p_{\bot}%
}e_{b0,\Vert}J_{0}\left(  \frac{k_{b\bot}v_{\bot}}{\Omega}\right)
\end{equation}
and, under the asumption that the lower hybrid wave is mostly electrostatic
(), we are left with
\begin{equation}
\Theta_{\mathbf{k}}^{b,\left(  0\right)  }\simeq\frac{p_{\Vert}}{p_{\bot}%
}e_{b0,\Vert}J_{0}\left(  \frac{k_{b\bot}v_{\bot}}{\Omega}\right)
\end{equation}


Note that we can modify the expression of the resonance condition as such:%
\begin{align*}
\frac{c}{\omega_{b}}\frac{\gamma m_{e}}{p\left\vert \xi\right\vert }%
\delta\left(  N_{b\Vert}-N_{\Vert\text{res}}\right)   &  =\delta\left(
\frac{\omega_{b}}{c}v_{\Vert}N_{b\Vert}-\frac{\omega_{b}}{c}v_{\Vert}%
N_{\Vert\text{res}}\right) \\
&  =\delta\left(  \omega_{b}-v_{\Vert}k_{b\Vert}\right) \\
&  =\frac{1}{\left\vert k_{b\Vert}\right\vert }\delta\left(  \frac{\omega_{b}%
}{k_{b\Vert}}-\frac{p_{\Vert}}{\gamma m_{e}}\right) \\
&  =\frac{1}{v_{Te}\left\vert k_{b\Vert}\right\vert }\delta\left(
\frac{p_{\Vert}}{\gamma p_{Te}}-\frac{c}{v_{Te}N_{b\Vert}}\right)
\end{align*}


Using a normalized Gaussian to smoothen the delta function%
\[
\delta\left(  \frac{p_{\Vert}}{\gamma p_{Te}}-\frac{c}{v_{Te}N_{b\Vert}%
}\right)  \rightarrow\frac{1}{\sqrt{\pi}\Delta_{\text{lis}}}\exp\left[
-\frac{1}{\Delta_{\text{lis}}^{2}}\left(  \frac{p_{\Vert}}{\gamma p_{Te}%
}-\frac{c}{v_{Te}N_{b\Vert}}\right)  ^{2}\right]
\]
we have finally the following transformation, to relate to Delphine:%
\[
\frac{\gamma m_{e}}{p\left\vert \xi\right\vert }\delta\left(  N_{b\Vert
}-N_{\Vert\text{res}}\right)  =\frac{\omega_{b}}{c}\frac{1}{v_{Te}\left\vert
k_{b\Vert}\right\vert }\frac{1}{\sqrt{\pi}\Delta_{\text{lis}}}\exp\left[
-\frac{1}{\Delta_{\text{lis}}^{2}}\left(  \frac{p_{\Vert}}{\gamma p_{Te}%
}-\frac{c}{v_{Te}N_{b\Vert}}\right)  ^{2}\right]
\]


Putting all terms together in (), we find%
\begin{align}
\frac{\overline{D}_{b,0}^{\text{RF}}(p,\xi)}{2\pi\delta\left(  \theta
-\theta_{b}\right)  }  &  =\frac{\omega_{b}}{c}\frac{1}{4\pi^{2}rR}\left\vert
\widehat{\psi}\cdot\widehat{r}\right\vert \frac{f_{\text{inc},b}^{l+1/2}%
}{\left\vert \mathbf{\Phi}_{b}\right\vert }P_{b,\text{inc}}\left\vert
\frac{p_{\Vert}}{p_{\bot}}e_{b0,\Vert}\right\vert ^{2}\\
&  \times\left[  J_{0}\left(  N_{b\bot}\frac{v_{Te}}{c}\frac{\omega_{b}%
}{\omega_{ce}}\frac{p_{\bot}}{p_{Te}}\right)  \right]  ^{2}\frac{\sqrt{\pi}%
}{\Delta_{\text{lis}}}\frac{e^{2}}{\omega_{b}\epsilon_{0}m_{e}^{2}\nu
_{e}v_{Te}^{3}}\frac{1}{\left\vert k_{b\Vert}\right\vert }\exp\left[
-\frac{1}{\Delta_{\text{lis}}^{2}}\left(  \frac{p_{\Vert}}{\gamma p_{Te}%
}-\frac{c}{v_{Te}N_{b\Vert}}\right)  ^{2}\right]
\end{align}


Now, if we consider circular concentric flux surfaces, we have that the metric
factor $\left\vert \widehat{\psi}\cdot\widehat{r}\right\vert =1$ and that the
incremental volume of the flux-surface under consideration is%
\[
\Delta V=4\pi^{2}rR\Delta r
\]
so that we can rewrite%
\begin{equation}
\frac{\overline{D}_{b,0}^{\text{RF}}(p,\xi)}{2\pi\delta\left(  \theta
-\theta_{b}\right)  }\frac{p_{\bot}^{2}}{p_{\Vert}^{2}}=\frac{2P_{b,\text{inc}%
}\Delta t}{\Delta V}\tau^{-1}\left[  J_{0}\left(  N_{b\bot}\frac{v_{Te}}%
{c}\frac{\omega_{b}}{\omega_{ce}}\frac{p_{\bot}}{p_{Te}}\right)  \right]
^{2}\frac{\sqrt{\pi}}{\Delta_{\text{lis}}}\frac{e^{2}}{\omega_{b}\epsilon
_{0}m_{e}^{2}\nu_{e}v_{Te}^{3}}\frac{1}{\left\vert k_{b\Vert}\right\vert }%
\exp\left[  -\frac{1}{\Delta_{\text{lis}}^{2}}\left(  \frac{p_{\Vert}}{\gamma
p_{Te}}-\frac{c}{v_{Te}N_{b\Vert}}\right)  ^{2}\right]
\end{equation}
with%
\[
\tau^{-1}=\frac{\Delta r}{2\Delta t}\frac{\omega_{b}}{c}\frac{f_{\text{inc}%
,b}^{l+1/2}}{\left\vert \mathbf{\Phi}_{b}\right\vert }\left\vert e_{b0,\Vert
}\right\vert ^{2}%
\]
while the term calculated by Delphine is
\[
\tau_{D}^{-1}=\frac{\left[  \left(  K_{\bot}-N_{b\Vert}^{2}\right)  \left(
K_{\bot}-N_{b}^{2}\right)  -K_{xy}^{2}\right]  }{\partial D/\partial\omega}%
\]


List of approximations made in the calculation of the quasilinear diffusion
coefficient in Delphine:

\begin{itemize}
\item Circular concentric flux surfaces.

\item Only the parallel polarization is accounted for (equivalent to making
the ES approximation). This is OK.

\item Only the first order thermal correction are included (to be verified,
probably OK).

\item The poloidal location of the ray is NOT accounted for in the resonance
condition (the approximation made here is that the ray power is uniformly
distributed on the flux-surface). This approximation is probably NOT\ OK.
\end{itemize}

\subsection{\bigskip Calculation of the power absorbed along the ray}

List of approximations made in the calculation of the absorption along the ray

\begin{itemize}
\item All approximations made in the DQL coefficient.

\item Small FLR effects (probably OK)

\item only the parallel distribution is taken into account

\item only the absorption at $N_{b\Vert}-N_{\Vert\text{res}}$ is accounted
for. (not the $\Delta_{\text{lis}}$ effect)

\item apparently, no appropriate flux-surface averaging
\end{itemize}


\end{document}
